% Standalone circuitikz version of potentiometer-multimeter_schem, for
% comparison against the schemdraw (../potentiometer_multimeter_schem.py)
% version. Build with build.sh in this directory.
%
% The meter is a plain hand-drawn circle+label, NOT circuitikz's built-in
% `ohmmeter` bipole: that bipole always draws a diagonal needle/arrow
% across the circle (baked into \drawmeteringcircle in the circuitikz
% source — it's an analog-gauge symbol, not a plain "Ω" badge), which
% doesn't match this site's other multimeter diagrams. A plain node gives
% full control and avoids it entirely.
%
% Also avoids a picture-level `scale=` key: combining `scale=` with
% circuitikz's bipoles here caused the fixed-size meter symbol and the
% scaled coordinate grid to disagree about where the connecting wires
% should land, producing visibly diagonal "funnel" wires into the meter
% instead of clean right angles. Sizing is done directly in coordinates
% and `\Large`/font choice instead.
\documentclass[tikz,border=4pt]{standalone}
\usepackage{circuitikz}
\usepackage[scaled]{helvet}
\renewcommand{\familydefault}{\sfdefault}

\definecolor{navy}{HTML}{2C3E50}

% Match the zigzag's line width to plain wires: circuitikz draws bipole
% bodies (resistor/potentiometer zigzags) at `bipoles/thickness` times the
% ambient line width (default thickness=2, i.e. twice as thick as a plain
% \draw at the same line width) — set it to 1 so the zigzag isn't
% noticeably thicker than the wires connecting to it.
\ctikzset{bipoles/thickness=1}

\begin{document}
\begin{circuitikz}[american, color=navy, thick, line width=1pt]
    \large

    % potentiometer: left leg A, right leg B, wiper tap
    \coordinate (A) at (0,0);
    \coordinate (B) at (4,0);
    \draw (A) to[potentiometer, name=POT] (B);

    % meter sits directly above A and the wiper (same x's as the legs it
    % connects to), so the connecting wires below are plain verticals
    \coordinate (M1) at (0,2.8);
    % (A -| B) = (B's x, A's y) -- the x and y each come from the
    % *other* operand than you'd guess from reading left-to-right.
    \coordinate (M2) at (M1 -| POT.wiper);
    \node[draw, circle, minimum size=1.3cm, line width=1pt] (OHM) at ($(M1)!0.5!(M2)$) {$\Omega$};

    \draw (A) -- (M1) -- (OHM.west);
    \draw (POT.wiper) -- (M2) -- (OHM.east);

    % junction dots (only where a real connection exists)
    \draw (A) node[circ]{};
    \draw (POT.wiper) node[circ]{};

    % labels
    \node[above] at (OHM.north) {\shortstack{Multimeter\\(resistance mode)}};
    \node[below] at (POT.south) {\shortstack{potentiometer\\10 k$\Omega$}};
    \node[below left] at (A) {leg 1};
    \node[below right] at (B) {\shortstack[l]{leg 2\\(not connected)}};
    \node[right] at ($(POT.wiper)+(0.15,0.25)$) {wiper (middle leg)};

\end{circuitikz}
\end{document}
